\documentclass[12pt,a4paper,oneside]{ctexart}
\usepackage{amsmath,amsthm,amssymb,bm,color,framed,graphicx,hyperref,mathrsfs,fancyhdr}

\title{\textbf{算法基础 Homework 5}}
\author{郝思粤 \\ 学号 PB23151779 \\ 中国科学技术大学}
\date{\today}
\linespread{1.5}
\definecolor{shadecolor}{RGB}{241,241,255}

% 页眉页脚
\pagestyle{fancy}
\fancyhf{}
\fancyhead[L]{算法基础 Homework 5}
\fancyhead[R]{PB23151779}
\fancyfoot[C]{\thepage}

% Problem 环境
\newcounter{problemname}
\newenvironment{problem}{
    \begin{shaded}
    \stepcounter{problemname}
    \addcontentsline{toc}{section}{题目\arabic{problemname}}
    \noindent\textbf{题目\arabic{problemname}. }
}{
    \end{shaded}\par
}
\newenvironment{solution}{
    \addcontentsline{toc}{subsection}{解答}
    \noindent\textbf{解答. }
}{\par}
\newenvironment{note}{
    \addcontentsline{toc}{subsection}{注记}
    \noindent\textbf{题目\arabic{problemname}的注记. }
}{\par}

\begin{document}
\maketitle
%\tableofcontents
\newpage

\maketitle

\section*{Q1.（2 × 10′ = 20 分）}

\subsection*{(a) 求矩阵链最优括号化方案（给出计算过程）}

给定维度序列
\[
p=(5,10,3,12,5,50,6),
\]
对应 6 个矩阵：
\[
A_1:5\times 10,\
A_2:10\times 3,\
A_3:3\times 12,\
A_4:12\times 5,\
A_5:5\times 50,\
A_6:50\times 6.
\]

令
\[
m[i,j]=\text{计算 }A_i A_{i+1}\cdots A_j\text{ 的最小标量乘法次数},\qquad m[i,i]=0.
\]
递推式
\[
m[i,j]=\min_{i\le k<j}\left(m[i,k]+m[k+1,j]+p_{i-1}p_kp_j\right).
\]

下面按链长 $L$ 计算。

\subsubsection*{链长 $L=2$}
\[
\begin{aligned}
m[1,2]&=5\cdot 10\cdot 3=150,\\
m[2,3]&=10\cdot 3\cdot 12=360,\\
m[3,4]&=3\cdot 12\cdot 5=180,\\
m[4,5]&=12\cdot 5\cdot 50=3000,\\
m[5,6]&=5\cdot 50\cdot 6=1500.
\end{aligned}
\]

\subsubsection*{链长 $L=3$}
\[
\begin{aligned}
m[1,3]&=\min\{0+360+5\cdot10\cdot12,\ 150+0+5\cdot3\cdot12\}=330,\\
m[2,4]&=330,\quad
m[3,5]=930,\quad
m[4,6]=1860.
\end{aligned}
\]

\subsubsection*{链长 $L=4$}
\[
m[1,4]=405,\qquad
m[2,5]=2430,\qquad
m[3,6]=1770.
\]

\subsubsection*{链长 $L=5$}
\[
m[1,5]=1655,\qquad
m[2,6]=1950.
\]

\subsubsection*{链长 $L=6$（完整矩阵链）}
\[
\begin{aligned}
m[1,6]&=\min\{
0+1950+5\cdot10\cdot6,\;
150+1770+5\cdot3\cdot6,\;
330+1860+5\cdot12\cdot6,\\
&\qquad\qquad 405+1500+5\cdot5\cdot6,\;
1655+0+5\cdot50\cdot6
\}\\
&=\boxed{2010}.
\end{aligned}
\]

故最优值 $m[1,6]=2010$。

回溯割点，可得到括号结构：
\[
(A_1A_2)\big((A_3A_4)(A_5A_6)\big),
\]
总乘法次数 $2010$。

\subsection*{(b) 证明：对 $n$ 个元素的表达式完全括号化，恰好需要 $n-1$ 对括号}

用归纳法。

\textbf{基例}：$n=1$ 时只有一个元素，不需要括号，$0=n-1$ 成立。

\textbf{归纳假设}：对任意 $k<n$，含 $k$ 个元素的完全括号化需要 $k-1$ 对括号。

\textbf{归纳步骤}：考虑含 $n$ 个元素的完全括号化，最外层一定是
\[
(E_L E_R),
\]
设左、右表达式分别含 $k$ 和 $n-k$ 个元素。根据假设，$E_L$ 需要 $k-1$ 对括号，
$E_R$ 需要 $(n-k)-1$ 对括号，外层再加一对，总数
\[
(k-1)+(n-k-1)+1=n-1.
\]
因此对 $n$ 也成立，命题得证。

\section*{Q2（20 分）}

参考提示为，设
\[
f(k,m)=\text{在最多 $m$ 次实验、$k$ 个蛋的情况下，最多可确定的楼层数}.
\]

一次实验有两种结果：

\begin{itemize}
\item 蛋碎：变成 $k-1$ 个蛋、$m-1$ 次实验，能确定 $f(k-1,m-1)$ 层；
\item 不碎：还是 $k$ 个蛋、$m-1$ 次实验，能确定 $f(k,m-1)$ 层。
\end{itemize}

再加上当前这一层，有
\[
f(k,m)=f(k-1,m-1)+f(k,m-1)+1.
\]

边界：
\[
f(0,m)=0,\quad f(k,0)=0.
\]

\textbf{求最少实验次数：} 从 $m=0,1,2,\dots$ 依次往上，对每个 $m$ 用上式滚动计算
$f(1,m),\dots,f(k,m)$，直到 $f(k,m)\ge n$ 为止，此时的 $m$ 即答案。每个 $m$ 只算
$k$ 个状态，所以总复杂度 $O(km^\ast)$，其中 $m^\ast$ 为最终答案。

进一步可以证明
\[
f(k,m)=\sum_{i=1}^k \binom{m}{i},
\]
它是关于 $m$ 的 $k$ 次多项式，所以大致有 $m^\ast=\Theta(n^{1/k})$，从而时间
$O(km^\ast)=O(n^{1/k})$，确实小于 $O(n)$。

\section*{Q3（20 分）}

给定长度为 $n$ 的数组 $\textit{nums}[1..n]$，每个数在 $[1,n]$，
要求找一个最长子序列，使得：

\begin{itemize}
    \item 严格递增；
    \item 相邻元素差不超过 $k$。
\end{itemize}

\textbf{状态：}
\[
dp[i]=\text{以第 $i$ 个元素为结尾的、满足条件的最长子序列长度}.
\]
初值 $dp[i]=1$。

\textbf{转移：} 对每个 $i$，枚举 $j<i$，若
\[
\textit{nums}[j]<\textit{nums}[i],\quad
\textit{nums}[i]-\textit{nums}[j]\le k,
\]
则可以接在后面：
\[
dp[i]=\max(dp[i],\,dp[j]+1).
\]

最后答案为 $\max_i dp[i]$。双重循环，时间复杂度 $O(n^2)$，空间 $O(n)$。

\section*{Q4（20 分）叠叠乐}

\subsection*{(a)}

把每本书看成一个二元组 $(a_i,b_i)$，$a_i$ 是长边，$b_i$ 是短边，且 $a_i>b_i$。
若书 $i$ 放在书 $j$ 上面，需要
\[
a_i<a_j,\quad b_i<b_j.
\]

为方便处理，先把所有书放到数组 $A[1..n]$ 中，然后按
\[
a\text{ 降序},\quad a\text{ 相同时 }b\text{ 降序}
\]
排序。排序后有
\[
A[1].a\ge A[2].a\ge\cdots\ge A[n].a,
\]
且同一长边里短边也是非增的。

接着做 DP。设
\[
C[i]=\text{以第 $i$ 本书为最上层时，最多能叠的层数},
\]
初值 $C[i]=1$。对每个 $i$，枚举 $j<i$，若
\[
A[i].a<A[j].a,\quad A[i].b<A[j].b,
\]
则可以把 $i$ 放到 $j$ 上面：
\[
C[i]=\max(C[i],\,C[j]+1).
\]

最后结果为 $\max_i C[i]$。排序 $O(n\log n)$，DP 双循环 $O(n^2)$，总复杂度为 $O(n^2)$。

\subsection*{(b)}

第 $i$ 种积木尺寸为 $a_i\times b_i\times c_i$，可以任意旋转。对每一种尺寸，我取出三块，
固定成三种摆放方式：

\[
a_i\times b_i\ (\text{高 }c_i),\quad
a_i\times c_i\ (\text{高 }b_i),\quad
b_i\times c_i\ (\text{高 }a_i).
\]

对每种摆放方式，把底面两条边从大到小排序成 $(\ell,w)$，$\ell\ge w$，这样一共
得到至多 $3n$ 个“虚拟积木”，每一个都有一个底面 $(\ell,w)$。

上层积木的底面必须在两个方向都严格小于下层底面，即
\[
(\ell_2,w_2)\text{ 可以放在 }(\ell_1,w_1)\text{ 上} \iff
\ell_2<\ell_1,\ w_2<w_1.
\]

这就和 (a) 完全一样，只是物品个数变成 $3n$。因此对这 $3n$ 个二元组 $(\ell,w)$ 套用
(a) 中的排序 + DP，得到的最长长度就是最多能叠的积木数。时间复杂度仍是
$O((3n)^2)=O(n^2)$。

\section*{Q5（2 × 10′ = 20 分）}

\subsection*{(a) 一维最大连续子数组}

设
\[
f(i)=\text{所有从位置 $i$ 开始的连续子数组中，和的最大值},\quad 1\le i\le n.
\]

考虑第 $i$ 个位置时，要么单独从这里重新开始，要么把后面的最优段接上：
\[
f(i)=
\begin{cases}
A[i], & f(i+1)<0,\\[4pt]
A[i]+f(i+1], & f(i+1)\ge 0.
\end{cases}
\]
边界 $f(n)=A[n]$。从 $i=n-1$ 一直算到 $1$，即可得到所有 $f(i)$，答案是
\[
\max_{1\le i\le n} f(i).
\]

如果想恢复具体区间，可以记录取得最大值的起点 $i^*$，然后从 $i^*$ 开始往右累加，
直到再加会变小为止。整体时间 $O(n)$。

\subsection*{(b) 二维最大子矩阵}

矩阵为 $M[1..n][1..n]$。套路是“先压成一维，再用上一问的算法”。

\textbf{第一步：行前缀和。}  
对每一行做前缀和
\[
S[i][j]=\sum_{t=1}^j M[i][t],\quad 1\le i,j\le n,
\]
这样任意一行在列区间 $[c,d]$ 上的和可以 $O(1)$ 算出：
\[
\sum_{t=c}^d M[i][t]=S[i][d]-S[i][c-1].
\]

\textbf{第二步：枚举左右列，把二维压成一维。}  
用两重循环枚举所有列区间 $[c,d]$。对固定一对 $(c,d)$，对每一行 $k$ 定义
\[
C[k]=S[k][d]-S[k][c-1],\quad 1\le k\le n.
\]
这相当于把列区间 $[c,d]$ 里每一行的和压成了一个一维数组 $C[1..n]$。

\textbf{第三步：在 $C$ 上跑 (a)。}  
在数组 $C$ 上用 (a) 中的 $f(i)$ 算法求最大连续子数组和（顺便可以记录对应的起止行）。
这个子数组对应的就是当前列区间 $[c,d]$ 下的最佳子矩阵。对所有 $(c,d)$ 取最大即可。

\textbf{复杂度：}  
前缀和 $O(n^2)$；列区间有 $O(n^2)$ 种，每种需要 $O(n)$ 时间构造 $C$ 并做一次一维 DP，
所以总时间是
\[
O(n^2)\cdot O(n)=O(n^3).
\]

\end{document}
